% Replacement text for Sections 6 and 7. Do not add a second document environment.
% Required packages in the main preamble: amsmath, amssymb, booktabs, graphicx.
% Keep the existing bibliography entries: tpch, tpcds, job, sqlstorm25.
% Figure paths are relative to the main compilation directory, not this file.
% Override this command before \input if the figures are copied to another project.
\providecommand{\DuckOuterYanFigureDir}{auto_test/results/plots/20260908_235240/t32}
% Data: experiment_onoff_benchmark_20260908_235240_t32.csv + .attempts.csv.
% The run identifier contains 0908; the t32 log records a 2026-09-10 start.
% Phase data: experiment_onoff_benchmark_20260830_002656.csv, 48 threads.
% Implementation checked against outer_yan_optimizer.hpp/.cpp and outer_yan_dp.cpp.

\section{Optimization}
\label{sec:optimization}

OuterYan compares two alternatives: DuckDB's normally optimized plan and a
reordered plan with semi-join reduction. It first estimates the reduction,
then chooses a join order for the reduced inputs, and finally checks whether
the expected join savings justify the reduction overhead.

\subsection{Cost Model}
\label{sec:cost-model}

We score a join tree by the sum of its estimated intermediate cardinalities:
\begin{equation}
 C_{\mathrm{out}}(T)=\sum_{\text{join }v\in T}\widehat{N}(v).
 \label{eq:cout}
\end{equation}
A leaf contributes zero cost. For an outer join, the output includes both
matching tuples and unmatched tuples from its preserved sides. The default
model assigns unit weight to each output tuple. These costs are row-work
estimates, not predictions of elapsed seconds.

Let $T_D$ denote the normally optimized DuckDB plan and $T_Y$ the OuterYan
candidate. We compare
\begin{equation}
 \begin{aligned}
 C_D &= C_{\mathrm{out}}(T_D),\\
 C_Y &= C_{\mathrm{out}}(T_Y)+\lambda C_{\mathrm{SJ}},
 \end{aligned}
 \label{eq:plan-costs}
\end{equation}
where $C_{\mathrm{SJ}}$ accounts for both reduction rounds and $\lambda=1$
in the evaluated configuration. Both plans use the same scoring function;
the baseline is not the unoptimized input tree. The reduction schedule is
fixed before join enumeration, so its cost is added once to the complete
candidate rather than to every dynamic-programming state.

\subsection{Cardinality Estimation}
\label{sec:cardinality-estimation}

For a join with input sizes $N_L,N_R$, let $d_L,d_R$ be the estimated numbers
of distinct join keys, bounded between one and the corresponding nonempty
input size. Composite-key estimates multiply the per-column distinct counts
under an independence assumption. Statistics remain local to each original
join edge: we do not propagate transitive equalities across outer joins.

Assuming uniform key frequencies and containment of the smaller key domain
in the larger one, we estimate the matching output $M$ and unmatched tuples
$U_L,U_R$ as
\begin{equation}
 \begin{aligned}
 M   &= \frac{N_LN_R}{\max(d_L,d_R)},\\
 U_L &= N_L\max(0,1-d_R/d_L),\\
 U_R &= N_R\max(0,1-d_L/d_R).
 \end{aligned}
 \label{eq:join-cardinality}
\end{equation}
An inner join returns $M$; a left, right, or full outer join additionally
returns $U_L$, $U_R$, or both, respectively. Empty inputs are handled directly
using the join semantics. Swapping the children also swaps LEFT and RIGHT,
so the same logical input remains preserved.

\subsection{Reduction and Plan Enumeration}
\label{sec:two-round-reduction}
\label{sec:plan-enumeration}

Reduction follows preservation semantics. An inner edge permits filtering
in both directions. A one-sided outer edge permits only the preserved side
to filter the non-preserved side; a full outer edge permits neither.
Provably redundant reductions are omitted. The complete bottom-up round
precedes the top-down round, and each step uses the estimates updated by
earlier steps.

For a filter built on $b$ and applied to $p$, we scale each edge-local
distinct count by that relation's retained fraction, with a floor of one
for a nonempty input. Using these updated distinct counts, the estimated
probe size becomes
\begin{equation}
 N'_p=N_p\min(1,d_b/d_p).
 \label{eq:semijoin-update}
\end{equation}
An empty build removes all probe tuples; an empty probe remains empty.
With unit row-work weights, the reduction cost is
\begin{equation}
 C_{\mathrm{SJ}}
 =\sum_{\text{new filters}}3N_b
  +\sum_{\text{probe steps}}(N_p+N'_p).
 \label{eq:semijoin-cost}
\end{equation}
The three build terms represent materialization, filter construction, and
replay. Probe work includes one lookup per input tuple and emission of the
survivors. Reusing a filter does not charge its build work again.

The final reduced sizes initialize the enumeration leaves. After rewriting
the join expression into an associative form, dynamic programming combines
connected relation subsets and retains the cheapest tree for each subset.
A combination costs the two child plans plus the estimated output of their
join, as in Equation~\ref{eq:cout}. Edge types and predicate orientation are
preserved, and cross products are never introduced. Output-root preferences
are tried where applicable; if they prevent a complete plan, the search
retries without that preference. For 12 or more relations, or after the
10,000-pair search budget is exhausted, a greedy search repeatedly merges
the cheapest connected pair.

For an all-inner tree with reductions in both directions on every edge,
the estimated output of a connected subplan is bounded below by its largest
reduced leaf. This prevents the uniform-key estimator from predicting fewer
tuples than full reduction allows. The correction is not applied to general
outer-join trees.

\subsection{Candidate Selection}
\label{sec:plan-selection}

We first compare the candidates' join groupings using the sets of stable
base-relation instances below their join nodes. This comparison ignores
join type, child orientation, and semi-join wrappers. If the groupings are
identical, DuckDB's plan is retained. A missing candidate, invalid cost
comparison, or proof that the baseline eliminates all join work also keeps
the baseline.

For a genuinely different grouping, the default gate requires $C_Y\leq C_D$.
If the reductions retain less than 90\% of the estimated input tuples in
total, the stricter condition $C_Y\leq0.90C_D$ requires additional headroom.
When a table filter uses a generic selectivity estimate, both candidates
are also rescored without that estimated filter reduction, and the scenario
with the larger candidate-to-baseline cost ratio governs selection.
An additional guard rejects a candidate that retains at least 99\% of its
input when the largest estimated input is at least 1,000 times the smallest
unfiltered base relation with such an uncertain filter.

A narrowly scoped stretch gate permits $C_Y\leq1.20C_D$ only when the
candidate's real-join cost is lower, its summed real-join input cardinality
is at most 30\% of the baseline's, and the reductions retain at least 95\%
of the input tuples. This gate is unavailable when the generic-filter
check, amplification guard, or full-reduction lower-bound correction is
active. Thus, the relaxed threshold targets benefits from regrouping, not
candidates that rely on aggressive reduction estimates. These thresholds
are fixed implementation parameters in the evaluated V7 policy; they are
heuristics, not guarantees of a runtime improvement.

\section{Experiments}
\label{sec:experiments}

We evaluate \emph{DuckOuterYan}, our integration of OuterYan into DuckDB,
using 32 worker threads. We report the performance of enabling the optimizer
over the complete measured workload, then examine selected plans and their
execution phases.

\subsection{DuckOuterYan System Design}
\label{sec:system-design}

DuckOuterYan integrates plan generation, selection, and execution into
DuckDB's existing vectorized engine. The optimizer operates after filter
pushdown and builds both alternatives from the same logical input. One
copy follows DuckDB's native outer-join simplification and join-order
optimization. The other is converted into an operator-tree representation
that records relation identities, equality predicates, and preservation
directions.

The OuterYan path simplifies outer joins where null rejection permits,
records the legal two-round reduction schedule, and rewrites join types
into an associative form for enumeration. Recording reductions before this
rewrite keeps their semantics independent of the subsequently chosen
grouping. After the cost comparison in Section~\ref{sec:optimization}, a
selected candidate is resimplified and lowered back to DuckDB operators;
otherwise the native alternative continues through the normal optimizer.
The active path handles eligible acyclic equality-join queries. Unsupported
queries, including queries containing aggregation or DISTINCT, use the
native optimizer.

At execution time, semi-join build operators materialize their input and
construct reusable membership filters. Probe operators filter DuckDB
\texttt{DataChunk}s without adding output columns or changing surviving
tuple multiplicities. Eligible bounded integer keys use a bitmap; other
keys, or ranges exceeding the per-filter memory budget, use a hash filter.
Pipeline dependencies prevent a probe from starting before its filter is
finalized and preserve the bottom-up/top-down phase boundary. The remaining
joins and result production use DuckDB's execution engine. Profiling
attributes work to the two reduction rounds, subsequent joins, and other
operators separately.

\subsection{Outer-Join Benchmark}
\label{sec:outerjoin-benchmark}

We derive a controlled core from TPC-H~\cite{tpch}, TPC-DS~\cite{tpcds},
JOB~\cite{job}, and SQLStorm~\cite{sqlstorm25}. Its 34 traceable seeds comprise
5 TPC-H, 7 TPC-DS, 10 JOB, and 12 SQLStorm queries. Each seed supplies a
connected acyclic skeleton of 3--6 relation instances. We retain its join
graph and emit five deterministic input-tree variants: primary left-deep
and right-deep trees, two connected bushy alternatives where available,
and an alternate right-deep traversal. When a bushy split is unavailable,
the corresponding slot uses a mixed left-deep or secondary right-deep tree.
The full core contains 170 queries: 46 left-deep, 80 right-deep, and 44 bushy.

Primary and control variants use full outer joins; the 12 mixed alternatives
retain their declared join types. Every active edge has a plain-column
equality, and generated selections are null-rejecting predicates in the
global WHERE clause. Explicit parentheses fix the input grouping. These
are independently specified outer-join queries derived from source graphs,
not semantically equivalent rewrites of the original benchmark queries or
necessarily of one another.

A supplementary component contains 36 normalized synthetic queries:
9 industry-inspired patterns and 28 historical optimizer stress cases,
with one overlap. The latter were originally selected for favorable
performance and are identified as such in the manifest. They are not
production traces or an unbiased sample of application queries. This
component includes the graph family and additional queries over the
TPC-H, TPC-DS, and JOB schemas. SQLStorm queries are grouped by their
underlying schema, not reported as a fifth family.

The full union has 206 queries. Removing 21 previously recorded
multithreaded timeouts produced the 185-query input to the reported run.
That run completed 176 correct timing pairs and encountered 9 further
timeouts, all during the OFF correctness execution at the 1,800-second
limit: 2 graph bushy queries, 5 TPC-H Q9 star-graph variants, and 2 TPC-DS
Q93 chain-graph variants. These nine have since been removed from the
active execution manifest, leaving 144 core and 32 supplementary queries.
The original results and source queries remain available. The active set
loses graph bushy and TPC-H core star coverage, so it is not described as
fully balanced. No successfully timed regression was removed.

We distinguish \emph{input tree shape} (left-deep, right-deep, or bushy)
from \emph{join-graph topology} (chain, star, or branch). The plotting category
\emph{industry} denotes seven industry-inspired queries, not another tree
shape or all 32 supplementary queries.

\subsection{Experimental Setup}
\label{sec:experimental-setup}

\paragraph{Environment and data.}
The reported server configuration is an Intel Xeon Silver 4116 at
2.10\,GHz, with 24 physical cores, 48 hardware threads, and 192\,GB RAM.
We use the release build of the DuckDB-based research fork and set
\texttt{threads=32} for both configurations. TPC-H and TPC-DS use scale
factor 100. JOB uses the local IMDb/JOB database, and graph queries use
relations built from the SNAP Bitcoin Alpha edge list.
% VERIFY before submission: hardware is carried over from the supplied draft.
% The run log verifies release binary, threads=32, both SF=100, and graph=bitcoin.
% Supply OS, exact tested commit/version, compiler, storage/memory/spill settings,
% and JOB/graph snapshot identifiers from the server; do not infer them locally.

\paragraph{Compared configurations.}
\emph{DuckDB (OFF)} disables only the \texttt{outer\_yan} optimizer.
\emph{DuckOuterYan (ON)} enables it with the V7 gated policy described above;
the existing plots label this configuration \emph{OuterYan}.
An ON execution can still select DuckDB's native plan. Among the 176 valid
queries, 55 select OuterYan and 121 retain DuckDB. Of the 55 selections,
54 pass the strict gate and one uses the stretch gate. All 55 are recorded
as using bitmap filters; this run does not compare bitmap and hash kernels.

\paragraph{Correctness and timing.}
The harness executes each bare query under OFF and ON and compares the
sorted serialized result rows, preserving duplicate rows rather than
checking only cardinality. Sorting is external to the query; no CTAS,
COPY, or LIMIT wrapper changes the optimized plan. Only successful
comparisons proceed to timing.
Each configuration uses a fresh CLI process with one warm-up followed by
three measured executions. We report their median wall-clock time.
Queries run in trash-output mode, which fully consumes the result without
terminal formatting or result-file output. The timer covers query
processing, including optimization and result consumption, but excludes
CLI startup and session-setting commands. Profiling is disabled during
these measurements. No cold-cache reset or randomized OFF/ON ordering is
claimed.

\paragraph{Reporting.}
We reconcile the final CSV and attempts file by benchmark family and query
name: final records take precedence, and attempts supply otherwise missing
queries. In this run both contain the same 185 records, counted once.
Statistics use all 176 correct, finite, positive OFF/ON timing pairs,
including native-plan fallbacks and measured regressions. Timeouts are
reported separately, never substituted with the time limit. The selected-plan
subset is reported separately from both the full workload and the display
sample.

For each query, speedup is $S_q=T_{\mathrm{OFF},q}/T_{\mathrm{ON},q}$.
We summarize a group of $n$ queries using
\begin{equation}
 \begin{aligned}
 S_{\mathrm{geo}}&=\exp\!\left(\frac{1}{n}\sum_q\log S_q\right),\\
 S_{\mathrm{total}}&=\frac{\sum_q T_{\mathrm{OFF},q}}
                          {\sum_q T_{\mathrm{ON},q}}.
 \end{aligned}
 \label{eq:reported-speedups}
\end{equation}
The geometric mean gives every query equal weight in log space; the total-time
ratio captures the aggregate elapsed-time change over the same queries.
We also report arithmetic mean, median, and maximum speedup. These are
descriptive summaries, not confidence intervals or significance tests.

\subsection{Results}
\label{sec:experimental-results}

\subsubsection{Overall Performance}
\label{sec:overall-performance}

\begin{table*}[t]
\centering
\small
\setlength{\tabcolsep}{7pt}
\caption{32-thread OFF/ON speedups. The upper panel includes all valid
queries, including native-plan fallbacks; the lower panel includes only
queries selecting OuterYan. All speedup columns are multiplicative factors.}
\label{tab:runtime-by-benchmark}
\begin{tabular}{lrrrrrr}
\toprule
Family & $n$ & Geomean & Mean & Median & Max & Total-time ratio \\
\midrule
\multicolumn{7}{l}{\textit{All valid queries}} \\
Graph  &   5 & 3.3503 & 3.3509 & 3.3391 &  3.4550 & 3.3563 \\
TPC-H  &  48 & 0.9979 & 1.0006 & 0.9986 &  1.5082 & 0.9769 \\
TPC-DS &  41 & 1.3704 & 5.0319 & 1.0035 & 81.7573 & 1.4295 \\
JOB    &  82 & 1.1880 & 1.3292 & 1.0231 &  7.9904 & 1.5547 \\
All    & 176 & 1.2061 & 2.1595 & 1.0066 & 81.7573 & 1.3411 \\
\midrule
\multicolumn{7}{l}{\textit{OuterYan-selected queries}} \\
Graph  &  5 & 3.3503 &  3.3509 & 3.3391 &  3.4550 & 3.3563 \\
TPC-H  &  5 & 1.0258 &  1.0480 & 0.9677 &  1.5082 & 0.9387 \\
TPC-DS &  7 & 6.5998 & 24.6548 & 3.1821 & 81.7573 & 4.9506 \\
JOB    & 38 & 1.4663 &  1.7190 & 1.2504 &  7.9904 & 1.8602 \\
All    & 55 & 1.8530 &  4.7255 & 1.4691 & 81.7573 & 2.6059 \\
\bottomrule
\end{tabular}
\end{table*}

Table~\ref{tab:runtime-by-benchmark} separates optimizer-enabled performance
from selected-plan performance. Over all 176 valid queries, DuckOuterYan
achieves a geometric-mean speedup of $1.2061\times$ and a total-time ratio
of $1.3411\times$. The sum of query medians falls from 6,243.88 to
4,655.73 seconds, a 25.44\% reduction. There are 111 speedups and 65 slowdowns;
113 ratios lie within $[0.95,1.05]$. The median is close to one, whereas the
arithmetic mean is raised by a small number of large gains.

The 144 completed controlled-core queries alone yield $1.1877\times$
geometric-mean and $1.3044\times$ total-time speedups. The 32 supplementary
queries yield $1.2927\times$ and $2.0627\times$, respectively. This separation
matters because some supplementary sources were performance-selected.
The combined result characterizes this constructed workload, not the
unmodified standard benchmarks or an unbiased production workload.

\begin{figure*}[t]
\centering
\includegraphics[width=\textwidth]{\DuckOuterYanFigureDir/speedup_all_queries.pdf}
\caption{Speedup of every valid query at 32 threads, including native-plan
fallbacks. Colors and horizontal regions distinguish benchmark families;
subgroups distinguish input-tree shapes or the industry category. Points
are translucent, and the logarithmic vertical axis preserves both gains
and regressions. The dashed line denotes equal running time.}
\label{fig:speedup-all-queries}
\end{figure*}

Figure~\ref{fig:speedup-all-queries} shows the variation within each family.
The five graph queries improve consistently, while TPC-H is near parity by
geometric mean and slightly slower by total time. TPC-DS supplies most of
the absolute savings, approximately 1,278.12 seconds, despite a median
speedup near one. JOB saves 111.72 seconds. The 55 OuterYan-selected queries
achieve $1.8530\times$ geometric-mean and $2.6059\times$ total-time speedups,
but these conditional results do not replace the full-workload scores.

\begin{table}[t]
\centering
\small
\setlength{\tabcolsep}{3pt}
\caption{32-thread results by input category over all valid queries.
Selected is the number executing an OuterYan plan. Industry is not a tree
shape. Speedup columns are multiplicative factors.}
\label{tab:runtime-by-category}
\begin{tabular}{lrrrrr}
\toprule
Category & $n$ & Selected & Geomean & Median & Total \\
\midrule
Left-deep  & 45 & 13 & 1.2988 & 1.0062 & 2.1939 \\
Right-deep & 74 & 26 & 1.2406 & 1.0082 & 1.2112 \\
Bushy      & 50 & 15 & 1.1061 & 1.0056 & 1.0426 \\
Industry   &  7 &  1 & 1.0328 & 1.0151 & 1.0160 \\
\bottomrule
\end{tabular}
\end{table}

Table~\ref{tab:runtime-by-category} shows aggregate gains in all four
reporting categories, with the largest total-time improvement among
left-deep inputs. These groups differ in schema, graph topology, and data
selectivity, so their differences do not isolate a causal effect of tree
shape. In particular, the large TPC-DS gains occur in left-deep variants
over star-shaped join graphs, not in the separate graph benchmark family.

\subsubsection{Running Time Comparison}
\label{sec:runtime-comparison}

\begin{figure*}[t]
\centering
\includegraphics[width=\textwidth]{\DuckOuterYanFigureDir/runtime_outeryan_queries.pdf}
\caption{Unprofiled running times of DuckDB and DuckOuterYan at 32 threads
for 32 OuterYan-selected queries: 5 graph, 5 TPC-H, 7 TPC-DS, and 15 JOB.
Each pair of adjacent bars represents one query, with seconds on a
logarithmic axis. G/H/D/J identify the families; LD/RD/B identify input-tree
shape, and I identifies the industry category.}
\label{fig:runtime-outeryan-queries}
\end{figure*}

Figure~\ref{fig:runtime-outeryan-queries} targets up to 15 queries per family,
using only actual OuterYan selections with complete timing and matching
profile records. We distribute slots across the available categories and
choose evenly spaced ranks of OFF running time within each category, not
ranks of speedup or ON time. The first three families have only 5, 5, and
7 eligible selections, so all are shown; JOB contributes five queries per
tree shape. The 32-query display includes six regressions and is not the
population used to compute the tables.

The TPC-DS Q5 mixed variant D-LD-01 falls from 67.09 to 0.86 seconds
($77.9108\times$), and the Q80 mixed variant D-LD-03 falls from 654.17 to
8.00 seconds ($81.7573\times$). Their star graphs combine a selective fact
or return relation with several optional inputs. The large gains are
consistent with avoiding expensive intermediate joins, but the current
measurements do not separate the contributions of reduction and reordering.
JOB also benefits: J-RD-08 falls from 46.34 to 9.90 seconds, and J-B-08
from 3.58 to 1.12 seconds. Bushy grouping therefore does not inherently
prevent a gain.

TPC-H exposes a different outcome. Four of its five selected queries are
slower: three variants of SQLStorm 25509 and one of 17751. They add
31.00 seconds in total. H-RD-03, for example, rises from 103.06 to
121.68 seconds despite passing the estimated-cost gate. JOB has five
selected regressions, four bushy and one right-deep, adding 1.49 seconds
in total. The remaining 18 JOB slowdowns retain DuckDB's plan; they cannot
be attributed to execution of OuterYan's semi-joins. Additional optimizer
work, final-plan differences, and run-to-run variation require separate
checks. By absolute time, the TPC-H cases are the higher-priority target.

\subsubsection{Execution Time Details}
\label{sec:execution-time-details}

The current 32-thread run has no phase profiles. Figure~\ref{fig:phase-breakdown}
therefore retains the available August 30, 48-thread profiles as an
\emph{explicitly labeled placeholder}. It uses the same query names, short
IDs, and plotting order as Figure~\ref{fig:runtime-outeryan-queries}, but not
the same measurement run. Its baseline is the OFF median from that older
48-thread experiment, not the current 32-thread baseline. Matching is by
family and full query name; the archived results do not include SQL hashes
to verify byte-identical query text across the server runs.

\begin{figure*}[t]
\centering
\includegraphics[width=\textwidth]{\DuckOuterYanFigureDir/phase_breakdown_outeryan.pdf}
\caption{OuterYan execution time details for the same 32 query IDs as
Figure~\ref{fig:runtime-outeryan-queries}. This is a placeholder using
August 30 profiles at 48 threads, not current 32-thread measurements.
Round 1, Round 2, Round 3, and Other are accumulated operator CPU times.
Total Time is profiled wall latency; Baseline Time is the old experiment's
unprofiled 48-thread OFF median. CPU components must not be added to the
wall-time bars.}
\label{fig:phase-breakdown}
\end{figure*}

Round 1 and Round 2 contain the bottom-up and top-down semi-join build/probe
operators, respectively. Round 3 contains the subsequent real joins.
Other accounts for the remaining operator CPU time, including scans and
non-join processing. Their sum is accumulated CPU work across threads,
not end-to-end elapsed time. An absent round has no recorded operator work;
missing measurements are not replaced by zero. Total Time includes profiling
overhead and is distinct from the unprofiled ON median.

The older profiles distinguish several execution regimes. In the graph
sample, Round 1 has no operators and Round 2 contributes less than 0.01\%
of CPU work; real joins and other processing dominate. For D-LD-01, the
two reduction rounds together contribute about 35.9\%, while Round 3
contributes 12.8\%. D-LD-03 spends only 1.5\% in reduction and 15.3\% in
Round 3, with most CPU work assigned to Other. Thus, neither a large
reduction fraction nor a small join fraction alone determines net speedup.

For the TPC-H 25509 variants, reduction consumes approximately 17.5--18.3\%
of the old profiled CPU work and Round 3 consumes 61.2--63.1\%. Both the
preparatory work and the final joins remain substantial. J-B-09
(SQLStorm JOB 11914, bushy) spends 33.0\% in reduction and 61.7\% in
Round 3. In contrast, J-B-12 (the synthetic bushy4a variant) spends only
1.9\% in reduction, but 58.0\% in Round 3 and 40.0\% in Other. Its
multi-valued movie relationships produce approximately 41 million result
rows in the older run, motivating closer inspection of intermediate
cardinality and output-processing costs rather than only a larger
semi-join penalty.

These profiles motivate two distinct cost-model questions: whether
reduction overhead is repaid, and whether the chosen grouping still
creates expensive intermediate or final results. They do not establish
the cause of the current 32-thread regressions. Matching 32-thread profiles,
OFF/ON plans, and targeted ablations are needed to attribute those effects.

\paragraph{Summary.}
DuckOuterYan reduces aggregate elapsed time on the completed workload,
with concentrated gains in graph, TPC-DS, and JOB queries. Native-plan
fallbacks remain part of the evaluation, and selected-plan regressions
remain visible. The results support reduction-aware outer-join optimization
while identifying large-output joins and reduction overhead as separate
targets for improving plan selection.
